% BACKGROUND & RELATED WORK
% Last generated: 2026-01-15T05:34:34.872Z

\subsection{Decentralized Autonomous Organizations}

DAOs emerged from the intersection of blockchain technology and organizational theory. The concept was first articulated by \citet{buterin2014dao} as ``an entity that lives on the internet and exists autonomously, but also heavily relies on hiring individuals to perform certain tasks that the automaton itself cannot do.''

Early implementations like The DAO (2016) demonstrated both the potential and risks of on-chain governance, with its \$60M exploit highlighting the importance of rigorous mechanism design \citep{dhillon2017dao}. Modern DAOs have evolved significantly, with diverse governance structures including:

\begin{itemize}
    \item \textbf{Token-weighted voting}: Compound \citep{compound2020governance}, Uniswap \citep{uniswap2021governance}, Aave \citep{aave2020governance}
    \item \textbf{Optimistic governance}: Optimism \citep{optimism2022governance}, with bicameral structure
    \item \textbf{Conviction voting}: Giveth, 1Hive \citep{emmett2019conviction}
    \item \textbf{Holographic consensus}: DAOstack \citep{daostack2019holographic}
    \item \textbf{Dual governance}: Lido \citep{lido2022dualgovernance}, with staker veto rights
\end{itemize}

Recent empirical work has quantified voting power concentration in deployed DAOs \citep{fritsch2022analyzing, fritsch2024voting}. A comprehensive review by \citet{han2025review} synthesizes the emerging literature on DAO governance, highlighting the tension between decentralization ideals and observed centralization. \citet{cong2025centralized} provide empirical evidence that DAO governance often exhibits significant whale dominance, with large token holders able to unilaterally influence outcomes. Large-scale analyses of over 100 DAOs reveal that treasury assets have grown to \$24.5 billion while participation dynamics vary significantly by DAO category \citep{nabben2024future}.

\subsection{Multi-Agent Systems}

Our work builds on the multi-agent systems (MAS) literature, particularly work on emergent behavior in complex adaptive systems \citep{wooldridge2009introduction}. In MAS, autonomous agents interact according to local rules, producing system-level dynamics that may be difficult to predict analytically.

Key concepts from MAS relevant to DAO simulation include:
\begin{itemize}
    \item \textbf{Agent heterogeneity}: Agents differ in preferences, resources, and strategies
    \item \textbf{Bounded rationality}: Agents use heuristics rather than optimal computation
    \item \textbf{Emergent coordination}: System-level order arises from local interactions
    \item \textbf{Path dependence}: Historical trajectories constrain future possibilities
\end{itemize}

\subsection{Mechanism Design}

Mechanism design, or ``reverse game theory,'' studies how to design rules that achieve desired outcomes given strategic agent behavior \citep{nisan2007algorithmic}. Classical results include:

\begin{itemize}
    \item \textbf{Revelation Principle}: Any mechanism can be converted to a truthful direct mechanism
    \item \textbf{Gibbard-Satterthwaite}: No voting system satisfies strategy-proofness, non-dictatorship, and unrestricted domain \citep{gibbard1973manipulation, satterthwaite1975strategy}
    \item \textbf{VCG Mechanisms}: Achieving efficiency through appropriate transfers
\end{itemize}

DAO governance mechanisms must navigate these theoretical constraints while accommodating real-world considerations like gas costs, participation fatigue, and adversarial behavior. These challenges echo classic collective action problems \citep{olson1965logic} and commons governance dilemmas \citep{ostrom1990governing}.

\subsection{Computational Social Choice}

Social choice theory studies collective decision-making, with computational social choice extending this to algorithmic settings \citep{brandt2016handbook}. Relevant topics include:

\begin{itemize}
    \item \textbf{Voting rules}: Plurality, approval, ranked choice, quadratic \citep{weyl2018quadratic, lalley2018quadratic, buterin2019flexible, daian2024quadratic}
    \item \textbf{Aggregation}: Arrow's impossibility theorem and its implications \citep{arrow1950difficulty}
    \item \textbf{Manipulation}: Strategic voting, whale attacks, and collusion resistance \citep{barbereau2024whale}
    \item \textbf{Liquid democracy}: Delegative voting systems and their governance implications \citep{hardt2023liquid, sharma2024delegated}
\end{itemize}

\subsection{Agent-Based Modeling in Economics}

Agent-based computational economics (ACE) provides methodological foundations for our work \citep{tesfatsion2006handbook}. ACE models have been used to study:

\begin{itemize}
    \item Market dynamics and price formation
    \item Emergence of institutions and norms
    \item Policy evaluation under bounded rationality
    \item Network effects in economic systems
\end{itemize}

Our simulation framework extends this tradition to the specific domain of decentralized governance.

\subsection{Prior DAO Simulation Work}

Limited prior work has applied simulation to DAO governance:

\begin{itemize}
    \item \citet{faqir2021cadcad} used cadCAD for tokenomic modeling, focusing on token engineering and bonding curve dynamics
    \item \citet{tan2023agentdao} explored LLM-based agent DAOs, demonstrating emergent coordination in simple governance tasks
    \item DeepDAO \citep{deepdao2023analytics} provides empirical analytics across hundreds of DAOs, offering participation and treasury benchmarks
    \item Snapshot \citep{snapshot2021governance} off-chain voting data has enabled large-scale empirical studies of DAO participation patterns
    \item Various projects have modeled specific mechanisms (conviction voting, bonding curves) in isolation
\end{itemize}

Recent empirical studies have provided valuable benchmarks for simulation validation. \citet{kiayias2024democracy} conducted an empirical study of decentralized governance dynamics in the Internet Computer SNS ecosystem, finding participation rates between 50-70\% with generally increasing trends over time. \citet{rozas2024democratic} examined mechanisms for enhancing democratic participation in DAOs. The first European DAO Workshop \citep{dawo2025trends} identified four critical research streams: governance mechanisms, technical frameworks, value assessment, and legal dimensions.

Our framework's architecture draws on established ABM frameworks such as Mesa \citep{mesa2023framework}, adapted for the specific requirements of DAO governance simulation. Our work differs by providing a comprehensive, open-source framework supporting multiple mechanisms, agent heterogeneity, and systematic parameter exploration across the full governance lifecycle.
